% sections/sec1_intro.tex
\section{引言}
\label{sec:intro}

设 $R$ 为交换环，并设 $M$ 与 $N$ 为 $R$-模.（我们始终假定所考虑的环具有乘法幺元，且模均假定为幺模：对所有 $m \in M$ 均有 $1 \cdot m = m$.）直和 $M \oplus N$ 是模的一种加法运算.在此，我们引入一种乘法运算 $M \otimes_R N$，称为\newterm{tensor-product}{张量积}{tensor product}.我们首先描述模的张量积应当具有怎样的形式，严格的定义见第~\ref{sec:construction}~节.

张量积最早起源于向量空间，这也是物理学与工程学中唯一出现张量积的背景，因此我们首先叙述向量空间的张量积.设 $V$ 与 $W$ 为域 $K$ 上的向量空间，并分别选取 $V$ 的一组基 $\{e_i\}$ 与 $W$ 的一组基 $\{f_j\}$.张量积 $V \otimes_K W$ 定义为以形式符号 $e_i \otimes f_j$ 为基的 $K$-向量空间（我们在定义上宣告这些新符号是线性无关的）.因此，$V \otimes_K W$ 由形如 $\sum_{i,j} c_{ij} e_i \otimes f_j$ 的形式和组成，其中 $c_{ij} \in K$.$V \otimes_K W$ 中的元素称为\newterm{tensor}{张量}{tensor}.对于 $v \in V$ 与 $w \in W$，将 $v$ 与 $w$ 用 $V$ 和 $W$ 的原始基展开，然后将 $v \otimes w$ 展开，仿佛 $\otimes$ 是一种非交换乘法（且允许提出标量因子），以此方式得到的 $V \otimes_K W$ 中的元素定义为 $v \otimes w$.

例如，设 $V = W = \R^2 = \R e_1 + \R e_2$，其中 $\{e_1, e_2\}$ 是标准基.（对两个 $\R^2$ 我们使用相同的基.）则 $\R^2 \otimes_\R \R^2$ 是一个以 $e_1 \otimes e_1$、$e_1 \otimes e_2$、$e_2 \otimes e_1$ 与 $e_2 \otimes e_2$ 为基的 $4$ 维空间.若 $v = e_1 - e_2$ 且 $w = e_1 + 2e_2$，则
\begin{equation}
  v \otimes w = (e_1 - e_2) \otimes (e_1 + 2e_2) := e_1 \otimes e_1 + 2e_1 \otimes e_2 - e_2 \otimes e_1 - 2e_2 \otimes e_2. \label{eq:1.1}
\end{equation}

$v \otimes w$ 是否依赖于 $\R^2$ 基的选取？作为检验，另选一组基，例如 $e'_1 = e_1 + e_2$ 与 $e'_2 = 2e_1 - e_2$.则 $v$ 与 $w$ 可以表示为 $v = -\frac{1}{3}e'_1 + \frac{2}{3}e'_2$ 以及 $w = \frac{5}{3}e'_1 - \frac{1}{3}e'_2$.通过形式计算，
\begin{align*}
  v \otimes w &= \left(-\frac{1}{3}e'_1 + \frac{2}{3}e'_2\right) \otimes \left(\frac{5}{3}e'_1 - \frac{1}{3}e'_2\right) \\
  &= -\frac{5}{9}e'_1 \otimes e'_1 + \frac{1}{9}e'_1 \otimes e'_2 + \frac{10}{9}e'_2 \otimes e'_1 - \frac{2}{9}e'_2 \otimes e'_2,
\end{align*}
如果将这最后一个线性组合中 $e'_1$ 与 $e'_2$ 关于 $e_1$ 和 $e_2$ 的定义代入、完全展开并合并同类项，便会回到公式~\eqref{eq:1.1} 等号右侧的和式.这表明 $v \otimes w$ 在 $\R^2 \otimes_\R \R^2$ 中具有与基的选取无关的明确意义，尽管要一般性地证明这一点看似令人望而生畏.

在模的背景下，张量积可以类似向量空间的情形来描述，但 $\otimes$ 所应满足的性质必须在一般意义下规定，而不能仅仅依赖于基（因为模可能根本不存在基）：对于 $R$-模 $M$ 与 $N$，它们的张量积 $M \otimes_R N$（读作 ``$M$ tensor $N$'' 或 ``$M$ 与 $N$ 在 $R$ 上的张量积''）是一个由所有形如 $m \otimes n$（其中 $m \in M, n \in N$）的符号所生成的 $R$-模——这里不是作为一组基，而仅仅作为生成集\footnote{回忆一下，$R$-模的一个生成集是指其有限 $R$-线性组合能充满整个模的子集.生成集总是存在的，因为整个模本身就是一个生成集.}——并且这些符号满足分配律：
\begin{equation}
  (m + m') \otimes n = m \otimes n + m' \otimes n, \quad m \otimes (n + n') = m \otimes n + m \otimes n'. \label{eq:1.2}
\end{equation}
此外，乘以标量 $r \in R$ 可以放入 $\otimes$ 的任意一侧：对 $m \in M$ 与 $n \in N$，
\begin{equation}
  r(m \otimes n) = (rm) \otimes n = m \otimes (rn). \label{eq:1.3}
\end{equation}
因此，书写 $rm \otimes n$ 是毫无歧义的：它既表示 $r(m \otimes n)$，也表示 $(rm) \otimes n$.

应当将 $M \otimes_R N$ 中的公式~\eqref{eq:1.2} 与~\eqref{eq:1.3} 和直和 $M \oplus N$ 中的公式加以对比，在直和中：
\[
  (m + m', n) = (m, n) + (m', 0), \quad r(m, n) = (rm, rn).
\]
在 $M \oplus N$ 中，一个元素 $(m, n)$ 可分解为 $(m, 0) + (0, n)$，但 $M \otimes_R N$ 中的元素 $m \otimes n$ 却无法这样拆开.虽然 $M \oplus N$ 中的每个元素都是有序对 $(m, n)$，但在 $M \otimes_R N$ 中，通常存在远多于形如 $m \otimes n$ 的乘积元素.$M \otimes_R N$ 中的一般元素称为\termcn{张量}，它是有限个此类积的 $R$-线性组合\footnote{作为对比，考虑多项式环 $R[X, Y]$，其元素不仅包括乘积项 $f(X)g(Y)$，更包括这些乘积项的和式 $\sum_{i,j} a_{ij} X^i Y^j$.事实上，作为 $R$-模有 $R[X, Y] \cong R[X] \otimes_R R[Y]$（例~\ref{ex:poly-tensor}）.}：
\[
  r_1(m_1 \otimes n_1) + r_2(m_2 \otimes n_2) + \dots + r_k(m_k \otimes n_k),
\]
其中 $k \geq 1$，$r_i \in R$，$m_i \in M$ 且 $n_i \in N$.由于 $r_i(m_i \otimes n_i) = (r_i m_i) \otimes n_i$，我们可以将 $r_i m_i$ 重新记作 $m_i$，从而将上述线性组合改写为简单的和式：
\begin{equation}
  m_1 \otimes n_1 + m_2 \otimes n_2 + \dots + m_k \otimes n_k. \label{eq:1.4}
\end{equation}

在直和 $M \oplus N$ 中，相等性很容易定义：$(m, n) = (m', n')$ 当且仅当 $m = m'$ 且 $n = n'$.而在 $M \otimes_R N$ 中，形如~\eqref{eq:1.4} 的两个和式何时相等？在我们给出的张量积描述下，这并不容易判定，除了一种特殊情况：$M$ 与 $N$ 是自由 $R$-模，分别具有基 $\{e_i\}$ 与 $\{f_j\}$.在这种情况下，$M \otimes_R N$ 是以 $\{e_i \otimes f_j\}$ 为基的自由模，因此 $M \otimes_R N$ 中的每个元素都是（有限）和式 $\sum_{i,j} c_{ij} e_i \otimes f_j$（其中 $c_{ij} \in R$），且两个此类和式相等当且仅当同类项的系数相等.

为了在 $M$ 与 $N$ 没有基时描述 $M \otimes_R N$ 中的相等性，我们将利用 $M \otimes_R N$ 的\newterm{ump}{普遍映射性质}{universal mapping property}.张量积是代数学中第一个唯有通过普遍映射性质才能获得一致、合理理解的概念：$M \otimes_R N$ 是把 $M \times N$ 上的双线性映射转化为线性映射的万有对象.正如 Jeremy Kun \cite{kun} 所写，$M \otimes_R N$ 是所有从 $M \times N$ 出发的双线性映射的``守门人''（gatekeeper）.

在第~\ref{sec:bilinear}~节讨论双线性（与多重线性）映射之后，张量积的严格定义与构造将在第~\ref{sec:construction}~节中给出.张量积的具体实例见第~\ref{sec:examples}~节.在第~\ref{sec:properties}~节中，我们将展示张量积如何与模的其他构造相互作用.第~\ref{sec:base-ext}~节介绍\newterm{base-extension}{基变换}{base extension}（或称标量扩张）这一重要运算，它是利用张量积将 $R$-模转化为 $S$-模的过程，其中 $S$ 是另一个环.最后，在第~\ref{sec:physics}~节中，我们将介绍物理学中用于张量的指标记号.

以下是关于张量与张量积的简短历史.张量（tensor）一词源于拉丁语 \textit{tendere}，意为``拉伸''.1822 年，Cauchy 在连续介质力学中引入了柯西应力张量；1861 年，Riemann 在几何学中创立了黎曼曲率张量，但当时他们均未使用``张量''这一名称.1884 年，Gibbs \cite[第~3~章]{gibbs-vector} 引入了 $\R^3$ 中向量的张量积，称其为``不定积''（indeterminate product）\footnote{Gibbs 选用这一名称是因为用他的话来说，该乘积是``两个向量乘积的最一般形式''，因为它除了双线性之外不服从任何定律，而双线性是任何配称为乘积的运算所必须满足的.早在 1844 年，Grassmann 就创立了一种称为``开积''（open product）的特殊张量 \cite[第~3~章]{grassmann}.}，并将其应用于研究物体的应变.他在 1886 年将不定积推广到了 $n$ 维空间 \cite{gibbs-grassmann}.

Voigt 于 1898 年利用张量来描述晶体上的应力与应变 \cite{voigt}，``张量''一词在此处首次以其现代物理含义登场\footnote{用 $i, j, k$ 表示 $\R^3$ 的标准基，Gibbs 将形如 $a i \otimes i + b j \otimes j + c k \otimes k$（其中 $a, b, c > 0$）的和式称为正张量（right tensor）\cite[第~57~页]{gibbs-vector}，但我不知道这是否对 Voigt 的术语产生过影响.}.在几何学中，Ricci 在 19 世纪末使用了张量，他与 Levi-Civita 于 1901 年发表的论文 \cite{ricci-levicivita}（英译本见 \cite{ricci-english}）对 Einstein 创立广义相对论起到了关键作用.``张量''一词在物理学与数学中的广泛使用归功于 Einstein；而 Ricci 与 Levi-Civita 之前则用平淡的名称``体系''（systems）来称呼张量.Weyl 关于群论与量子力学的名著将向量空间 $V$ 与 $W$ 的张量积称为``乘积空间''（product space），并记作 $V \times W$ \cite[第~II~章第~10~节]{weyl-groups}.记号 $\otimes$ 归功于 Murray 与 von Neumann 在 1936 年关于 Hilbert 空间张量积（他们称之为``直积''）的研究 \cite[第~II~章]{murray-vonneumann}\footnote{感谢 Jim Casey 让我注意到文献 \cite{murray-vonneumann}.}.阿贝尔群 $A$ 与 $B$ 的张量积（使用了该名称，但记作 $A \circ B$ 而非 $A \otimes_\Z B$）由 Whitney 于 1938 年给出 \cite{whitney}.而交换环上模的张量积则由 Bourbaki 于 1948 年正式确立 \cite{bourbaki-alg}.
